%!TEX root = main.tex
%!TEX spellcheck = en_US

\begin{abstract}

Verifying network configurations is critical for ensuring operational stability, but traditional tools force a trade-off between the high manual effort of model checkers and the semantic shallowness of consistency checkers. While Large Language Models (LLMs) are a promising alternative, naive approaches like full-file, partition-based, or chain-of-thought prompting fail to handle the large, interconnected nature of configuration files, leading to context overload or fragmentation.
We introduce the Context-Aware Prompting (\sysname{}) framework, which enables an LLM to reason about network configurations like a human expert. \sysname{} first performs comprehensive context mining, identifying all potentially relevant neighboring, similar, and referenced (both intra- and cross-device) configurations. It then engages the LLM in a structured, interactive dialogue, where the model first requests the specific context it needs before performing its final, focused analysis. This on-demand approach provides the necessary context for deep reasoning while avoiding overload.
Our evaluation on production network configurations demonstrates that \sysname{} achieves near-perfect accuracy on known misconfigurations and uncovers over 20 previously undetected, real-world errors. We show that \sysname{} significantly outperforms existing tools and that the choice of LLM backend reveals a crucial precision-recall trade-off. Ultimately, \sysname{} provides a more effective and scalable paradigm for LLM-based network verification.

\end{abstract}
